\documentclass{article}

\usepackage{amsmath,amssymb,amsthm}
\usepackage{geometry}
\geometry{margin=1in}

\title{Tseitin Formulas on Expanders: Width Lower Bounds}
\author{URF Toolkit}
\date{\today}

\newtheorem{theorem}{Theorem}
\newtheorem{lemma}{Lemma}
\newtheorem{definition}{Definition}

\begin{document}
\maketitle

\section{Tseitin Contradictions on Expanders}

\begin{definition}
Let $G=(V,E)$ be a graph and $\chi:V\to\{0,1\}$ a charge function with
\[
\sum_{v\in V}\chi(v)\equiv 1 \pmod 2.
\]
The Tseitin formula $T(G,\chi)$ encodes parity constraints at each vertex.
\end{definition}

\section{Width Lower Bound Strategy}

\begin{theorem}[Target]
If $G$ is a constant-degree expander on $n$ vertices,
then any resolution refutation of $T(G,\chi)$
requires width
\[
W(T(G,\chi)) \ge \Omega(n).
\]
\end{theorem}

\section{Proof Skeleton}

\begin{lemma}[Expansion Implies Boundary Growth]
For any set of vertices $S\subseteq V$ with $|S|\le n/2$,
the edge boundary satisfies
\[
|\partial S| \ge \varepsilon |S|
\]
for some constant $\varepsilon>0$.
\end{lemma}

\begin{lemma}[Parity Preservation]
Any clause of width $w$ can constrain at most $w$ edge variables.
\end{lemma}

\begin{lemma}[Width Growth Mechanism]
To derive contradiction, a clause must encode a parity constraint
over a set of vertices whose induced boundary is empty.
Expansion forces such sets to be of size $\Omega(n)$.
\end{lemma}

\end{document}
